\section{The parameterized physical source}
\label{sec:v74-physical-graph}
We record explicitly the finite-dimensional family to which integration
and monotonicity are applied. The integer $m$ is fixed throughout this
section. A parameter uniformity assertion in this section is not an
assertion uniform in $m$.

\subsection{The moving section}
Let $a=(1,0)$, $b=(1/2,\sqrt3/2)$ and $\Lambda=\Z a+\Z b$.
In collision coordinates $(\alpha,p)$ the invariant probability is
$d\alpha\,dp/(4\pi)$. Put $h_0=1/200$. The four squares defining
$Y_R$ have half-side $h_0$ and centers
\[
 (0,0),\quad(\pi/6,-1/2),\quad
 (-\theta_R,\sin\theta_R),\quad
 (\pi/3-\theta_R,\sin\theta_R).
\]
The actual section is $Y_R^*=Y_R\cup U$, where
\[
 U=\{|\alpha-\pi/6|<3/50,\ |p|<3/20\},\qquad
 c_*:=\nu(Y_R^*)=\frac{91}{10000\pi}.
\]
Angles are taken modulo $2\pi$. These five rectangles are disjoint.
The probability used below is always $\nu_R^*=c_*^{-1}\nu|_{Y_R^*}$.

\begin{lemma}[An algebraic coordinate for the moving centers]
\label{lem:v74-half-angle}
For $R\in I=[9/20,47/100]$, the number $q_R=\tan(\theta_R/2)$
is the unique zero in $(49/100,53/100)$ of
\begin{equation}\label{eq:v74-quartic}
 P(R,q)=8Rq(1+q^2)-4q(1-q^2)
                       +\sqrt3(1-6q^2+q^4).
\end{equation}
In particular the normals and the $p$ coordinates of all five section
centers are semialgebraic functions of $R$.
\end{lemma}
\begin{proof}
The defining equation for the physical winding orbit is
$2R\sin\theta-\sin\theta\cos\theta
 +(\sqrt3/2)\cos(2\theta)=0$.
Substitute $\sin\theta=2q/(1+q^2)$ and
$\cos\theta=(1-q^2)/(1+q^2)$ and multiply by $2(1+q^2)^2$.
This gives \eqref{eq:v74-quartic}. On the displayed rectangle,
\[
 \partial_qP=8R(1+3q^2)-4+12q^2+\sqrt3(-12q+4q^3)<-2.
\]
For example, the positive terms are bounded above at $R=47/100$,
$q=53/100$, and the last term above by
$-(1732/1000)(12(49/100)-4(49/100)^3)$.
Also $P(R,49/100)>0$ and $P(R,53/100)<0$ throughout $I$.
Both endpoint signs follow by monotonicity in $R$ and the rational
bounds $1732/1000<\sqrt3<1733/1000$. Thus the zero exists and is
unique. The physical root, whose angle lies in $(0.92,0.967)$,
has its half-angle in this interval and hence is this zero.
The implicit derivative is nonzero, and the graph is semialgebraic.
The displayed rational sine and cosine, followed by fixed rotations,
give the asserted center coordinates.
\end{proof}

An angular interval of half-width $h<\pi$ and center normal $n_*$
can be tested by $n\cdot n_*>\cos h$. Its companion $p$ interval
uses two linear inequalities. Consequently the five rectangle tests
are semialgebraic in $R,n,p$. The fixed real constants $\cos h_0$
and $\cos(3/50)$ are coefficients; their presence does not require
an algebraic value of either angle.

\subsection{A finite graph with the exact labels}
Fix a uniform upper flight bound $\tau_+$ from the finite-horizon
geometry; the lower bound is $\tau_-=3/50$. Set $R_+=47/100$ and
\[
 \mathcal A=\{z\in\Lambda\setminus\{0\}:|z|\le\tau_++2R_++1\},
 \qquad D=\max_{z=z_1a+z_2b\in\mathcal A}\max(|z_1|,|z_2|).
\]
The set $\mathcal A$ is finite. It contains every next-center
increment and every center needed for first-hit and minimal
third-disk clearance tests. A center whose disk meets a segment
is within $\tau_++2R_+$ of its starting center. For the clearance,
at least one of the six nearest neighbors of that starting center
is not the terminal center. Its distance to the initial contact
is at most $1+R$. Thus the minimum third-center distance to the
segment is at most $1+R$. Any center attaining that minimum is
within $\tau_++2R_++1$ of the starting center. Centers outside
$\mathcal A$ cannot change the exact minimum. The starting disk
is tested separately for first-hit conditions.
Thus a finite set containing all attainable exact labels is
\begin{equation}\label{eq:v74-label-set}
 \mathscr L_m=\{(n,k):1\le n\le m,\ k\in\Z^2,
                                      \max(|k_1|,|k_2|)\le mD\}.
\end{equation}
Unattainable labels in this set give the zero source. No labels
are pooled in the integrations which follow.

Use a finite disjoint refinement of rational angular charts
\[
 n(\alpha)=\operatorname{Rot}_{\alpha_*}
       \left(\frac{1-x^2}{1+x^2},\frac{2x}{1+x^2}\right),
 \qquad |x|<1.
\]
Overlap boundaries are assigned by a fixed convention. For a center
word $C_0=0,C_1,\ldots,C_m$, keep as auxiliary variables the contact
positions $Q_j$, normals $n_j$, velocities $v_j$ and flight lengths
$\tau_j$. The graph includes
\begin{align*}
 Q_j&=C_j+Rn_j,& |n_j|^2&=|v_j|^2=1,\\
 Q_{j+1}&=Q_j+\tau_jv_j,& \tau_-&\le\tau_j\le\tau_+,\\
 v_j&=v_{j-1}-2(v_{j-1}\cdot n_j)n_j,
       &v_{j-1}\cdot n_j&<0\quad(1\le j\le m).
\end{align*}
At time zero, $v_0=c_0n_0+p\,n_0^\perp$, with
$c_0>0$ and $c_0^2+p^2=1$. Exclusion of an earlier collision is
imposed for every candidate disk. These are finite semialgebraic
conditions: the minimum distance to a segment is obtained at an
endpoint or at the orthogonal projection when that projection lies
inside the segment. Projection denominators are positive. The
intended endpoint disks retain the first-hit sign conditions; no
trajectory is continued through grazing. This is the same physical
first-hit graph as in the complete-height theorem.

The incidence margins are positive square roots or dot products on
this graph. The original clearance margins are the minima of the
corresponding finitely many candidate-distance branches, with their
branch conditions retained. Distances are encoded by a nonnegative
variable and its squared-distance equation. Hence the original
margins $c_j,\Delta_j$, not substitute margins, are semialgebraic.
The first-hit inequalities and regularity conditions remain part of
the domain when a margin tends to zero.

Write $\eta_j=\one_{Y_R^*}(T_R^jx)$. The exact event requires
\[
 \eta_0=\eta_m=1,\qquad
 \sum_{j=0}^{m-1}\eta_j=n,\qquad C_m=k_1a+k_2b.
\]
The terminal test is not an occupation summand. The physical roof
is $L=\sum_{j=0}^{m-1}\tau_j$. The graph projects to the two original
initial variables $(x,p)$; over a regular initial state the physical
successor is unique. Distinct physical center words therefore give
disjoint projected pieces, up to the already excluded singular
null set. Integration below is in these two initial variables,
not volume integration on the auxiliary graph.

\subsection{The exact integration family}
Put $N=2m+1$, $q_0=47/53$,
$e_j=q_0^{\min(j,m-j)/4}$ and
$e'_j=q_0^{\min(j,m-1-j)/4}$. The normalized margins are
$c_j/(\varepsilon e_j)$ and $\Delta_j/(\varepsilon e'_j)$.
Enumerate them as $a_i$, and for $y\in[1,2]^N$ set
\[
 W_y=1-\prod_{i=1}^N\one_{\{a_i>y_i\}}.
\]
Here the parameter domain is $R\in I$, $\varepsilon>0$,
$y\in[1,2]^N$, $t\in\R$, and $(n,k)\in\mathscr L_m$.
The physical application subsequently restricts $\varepsilon$ to
its original small range. Finiteness of this fixed-formula family
holds on the stated, larger positive domain.

\begin{theorem}[An explicit tame disintegration family]
\label{thm:v74-tame-family}
For each fixed $m$ and exact label, the sublevel mass
\begin{equation}\label{eq:v74-sublevel-mass}
 M(R,\varepsilon,y,t)
  =\int \one_{\Pi_{n,k,m,R}}W_y\one_{\{L\le t\}}\,d\nu_R^*
\end{equation}
is constructible. Its almost-everywhere roof derivative has a
representative definable in $\R_{\mathrm{an},\exp}$, bounded between
zero and $L_m=CA^m/c_*$. There is $J_m<\infty$, uniform over the
stated continuous parameters and the finite set \eqref{eq:v74-label-set},
which bounds the number of monotonicity intervals and point cells
of that representative. The complete original smooth-guard density
belongs to $BV(\R)$ and satisfies
\begin{equation}\label{eq:v74-bv}
 \|Db^\varepsilon\|_{\TV}
 \le4(J_m+1)L_m\|d\Theta\|_{\TV}^{2m+1}.
\end{equation}
\end{theorem}
\begin{proof}
On each disjoint initial chart the integrand supplied to the
parameterized integration theorem is exactly
\[
 \frac{\one_{\Pi_{n,k,m,R}}W_y\one_{\{L\le t\}}}
                         {2\pi c_*(1+x^2)},
\]
extended by zero off the projected regular physical graph.
All tests in this formula have just been specified. Projection
preserves semialgebraicity, and the deterministic graph makes the
roof and the tests single-valued. The integrand is globally
subanalytic and integrable in $(x,p)$ for every parameter.
The stability theorem for constructible functions under integration
\cite[Theorem 1.3]{CMInt} proves the assertion about $M$ after the
finite disjoint sum over charts and words. It does not assert that
$M$ itself is semialgebraic.

Positive domination by the complete density gives a Lipschitz
constant $L_m$ for each $t$ section. Define $g=\partial_tM$ where
the finite derivative exists, and $g=0$ elsewhere. Differentiability
and its derivative are definable by quantified inequalities in
$\R_{\mathrm{an},\exp}$. Thus $g$ is definable, and absolute
continuity identifies it with the positive threshold density almost
everywhere. Positivity and the Lipschitz bound give $0\le g\le L_m$.
The uniform monotonicity theorem for definable families
\cite{Coste} gives finitely many cells with the roof coordinate
last. Take the maximum over the finite labels. This defines $J_m$
for this fixed formula, without asserting an estimate as its
collision length grows.

Let $\sigma=d\Theta$. Although $\sigma$ may be signed, every $W_y$
is positive. The exact identity
$1-H^\varepsilon=\int W_y\,d\sigma^{\otimes N}(y)$ follows from
$\Theta(a)=\int\one_{\{y<a\}}\,d\sigma(y)$ and $\sigma(\R)=1$.
A bounded monotone-piece density extended by zero has variation
at most $4(J_m+1)L_m$. Integrating its distributional derivative
against the total variation of $\sigma^{\otimes N}$ proves
\eqref{eq:v74-bv} by the distributional characterization of $BV$
\cite{AFP}. This is the derivative of the full positive smooth-guard
source. Its support is contained in $[m\tau_-,m\tau_+]$, and its
total mass is zero. The support is compact at each fixed count;
it is not one compact interval independent of $m$.
\end{proof}

The theorem supplies the finite current needed below, including
incidence and clearance outside the selected disks. It supplies no
collision-normalized concentration estimate. In particular neither
$J_m$ nor the variation bound in \eqref{eq:v74-bv} is used as an
$m^{-2}$ estimate.
